\beforesec
\section{Background and experimental setup}\label{sec:expsetup}
\postsec

Our experiments compare the performance of zero-shot models, corresponding fine-tuned models, and models produced by WiSE-FT. 
To measure robustness, we contrast model accuracy on two related but different distributions, a reference distribution $\mathcal{D}_{\mathrm{ref}}$ which is the target for fine-tuning, and shifted distribution $\mathcal{D}_{\mathrm{shift}}$.\footnote{$\mathcal{D}_{\mathrm{ref}}$ and $\mathcal{D}_{\mathrm{shift}}$ are sometimes referred to as \textit{in-distribution} (ID) and \textit{out-of-distribution} (OOD). In this work, we include evaluations of zero-shot models, which are \textit{not} trained on data from the reference distribution, so referring to $\mathcal{D}_{\mathrm{ref}}$ would be imprecise. For clarity, we avoid the ID/OOD terminology.}
We assume both distributions have test sets for evaluation, and $\mathcal{D}_{\mathrm{ref}}$ has an associated training set $\mathcal{S}_{\mathrm{ref}}^\text{tr}$ which is typically used for training or fine-tuning.
The goal for a model is to achieve both high accuracy and consistent performance on the two distributions $\mathcal{D}_{\mathrm{ref}}$ and $\mathcal{D}_{\mathrm{shift}}$.
This is a natural goal as humans often achieve similar accuracy across the distribution shifts in our study \cite{shankar2020evaluating}.

For a model $f$, we let $\mathsf{Acc}_{\mathrm{ref}}(f)$ and $\mathsf{Acc}_{\mathrm{shift}}(f)$ refer to classification accuracy on the reference and shifted test sets, respectively. 
We consider $k$-way image classification, where $x_i$ is an image with corresponding label $y_i \in \{1,...,k\}$. 
The outputs of $f$ are $k$-dimensional vectors of non-normalized class scores.


\smallpara{Distribution shifts.}
\citet{taori2020measuring} categorized distribution shifts into two broad categories:
(i) \emph{synthetic}, e.g., $\ell_\infty$-adversarial examples or artificial changes in image contrast, brightness, etc. \cite{imagenetc,biggio2013evasion,biggio2018wild,geirhos2018generalisation,alcorn2019strike}; and 
(ii) \emph{natural}, where samples are not perturbed after acquisition and changes in data distributions arise through naturally occurring variations in lighting, geographic location, crowdsourcing process, image styles, etc. \cite{taori2020measuring, pmlr-v97-recht19a, imagenetr, imageneta, wilds2021}.
Following \citet{radford2019language}, our focus here is on natural distribution shifts as they are more representative of the real world when no active adversary is present. 
Specifically, we present our key results for five natural distribution shifts derived from ImageNet (i.e., $\mathcal{S}^\text{tr}_\text{ref}$ is ImageNet):
\begin{itemize}
 \item
ImageNet-V2 (IN-V2) \cite{pmlr-v97-recht19a}, a reproduction of the ImageNet test set with distribution shift
\item ImageNet-R (IN-R) \cite{imagenetr}, renditions (e.g., sculptures, paintings) for 200 ImageNet classes
\item ImageNet Sketch (IN-Sketch) \cite{imagenetsketch}, which contains sketches instead of natural images
\item ObjectNet \cite{objectnet}, a test set of objects in various scenes with 113 classes overlapping with ImageNet
\item ImageNet-A (IN-A) \cite{imageneta}, a test set of natural images misclassified by a ResNet-50 \cite{he2016deep} for 200 ImageNet classes.
\end{itemize}
Figure \ref{fig:data_samples} illustrates the five distribution shifts.

\smallpara{Effective robustness and scatter plots.}
To compare the robustness of models with different accuracies on the reference distribution, we follow the \emph{effective robustness} framework introduced by \citet{taori2020measuring}.
Effective robustness quantifies robustness as accuracy \emph{beyond a baseline} trained only on the reference distribution.
A useful tool for studying (effective) robustness are scatter plots that illustrate model performance under distribution shift  \cite{pmlr-v97-recht19a, taori2020measuring}. 
These scatter plots display accuracy on the reference distribution on the $x$-axis and accuracy under distribution shift on the $y$-axis, i.e., a model $f$ is shown as a point $\left(\mathsf{Acc}_{\mathrm{ref}}(f),\ \mathsf{Acc}_{\mathrm{shift}}(f)\right)$.
Figure \ref{fig:fig1} exemplifies these scatter plots with both schematics and real data.
For the distribution shifts we study, accuracy on the reference distribution is a reliable predictor of accuracy under distribution shift \cite{taori2020measuring, miller21b}. 
In other words, there exists a function $\beta : [0,1] \rightarrow [0,1]$ such that $\mathsf{Acc}_\mathrm{shift}\mleft(f\mright)$ approximately equals $\beta\mleft(\mathsf{Acc}_\mathrm{ref}\mleft(f\mright)\mright)$ for models $f$ trained on the train set $\mathcal{S}^\text{tr}_\mathrm{ref}$.  
Effective robustness \cite{taori2020measuring} is accuracy beyond this baseline, defined formally as $\rho(f) = \mathsf{Acc}_\mathrm{shift}(f) -\beta\mleft(\mathsf{Acc}_\mathrm{ref}\mleft(f\mright)\mright)$. 


In the corresponding scatter plots, effective robustness is vertical movement above expected accuracy under distribution shift (Figure \ref{fig:fig1}, top).
Effective robustness thereby disentangles accuracy changes on the reference distribution from the effect of robustness interventions.
When we say that a model is robust to distribution shift, we mean that effective robustness is positive. \citet{taori2020measuring} observed that no algorithmic robustness intervention consistently achieves substantial effective robustness across the distribution shifts in Figure~\ref{fig:data_samples}---the first method to do so was zero-shot CLIP.
Empirically, when applying logit (or probit) axis scaling, models trained on the reference distribution approximately lie on a linear trend \cite{taori2020measuring,miller21b}. As in \citet{taori2020measuring}, we apply logit axis scaling and show 95\% Clopper-Pearson confidence intervals for the accuracies of select points.



 \smallpara{Zero-shot models and CLIP.} We primarily explore CLIP models \cite{radford2021learning}, although we also investigate other zero-shot models including ALIGN \cite{jia2021scaling}, BASIC \cite{pham2021scaling} and a ViT model pre-trained on JFT \cite{dosovitskiy2021an}. Zero-shot models exhibit effective robustness and lie on a qualitatively different linear trend (Figure~\ref{fig:fig1}). 
CLIP-like models are pre-trained using image-caption pairs from the web.
Given a set of image-caption pairs $\{(x_1,s_1)...,(x_B,s_B)\}$, CLIP-like models train an image-encoder $g$ and text-encoder $h$ such that the similarity $\left\langle g(x_i), h(s_i)\right\rangle$ is maximized relative to unaligned pairs.
CLIP-like models perform zero-shot $k$-way classification given an image $x$ and class names $C = \{c_1,...,c_k\}$ by matching $x$ with potential captions.
For instance, using caption $s_i = \text{``a photo of a \{$c_i$\}''}$ for each class $i$, the zero-shot model predicts the class via $\argmax_j \left\langle g(x), h(s_j) \right\rangle$.\footnote{For improved accuracy, the embedding of a few candidate captions are averaged, e.g., $s_i^{(1)}{=}\text{``a \textit{photo} of a \{$c_i$\}''}$ and $s_i^{(2)}{=}\text{``a \textit{picture} of a \{$c_i$\}''}$ (referred to as prompt ensembling \cite{radford2021learning}).}
Equivalently, one can construct $\textbf{W}_{\text{zero-shot}} \in \mathbb{R}^{d \times k}$ with columns $h(s_j)$ and compute outputs $f(x) = g(x)^\top \textbf{W}_{\text{zero-shot}}$.
Unless explicitly mentioned, our experiments use the CLIP model \texttt{ViT-L/14@336px}, although all CLIP models are displayed in our scatter plots (additional details provided in Appendix~\ref{sec:morezs}).



